% ============================================================================
% SPANISCH LANGENTWURF - DS 7: PUFFER-STUNDE (CULTURA Y DIVERSIÓN)
% ============================================================================
% Kurs: Jahrgangsübergreifend Spanisch Spätbeginner (Kl. 10, 11, 12)
% Datum: Mo 26.01.2026
% Besonderheit: PUFFER-Stunde - entbehrlich für Beobachtungsstunde am 03.02.
% ============================================================================

\documentclass[11pt,a4paper]{article}

% ============================================================================
% PACKAGES
% ============================================================================
\usepackage[utf8]{inputenc}
\usepackage[T1]{fontenc}
\usepackage[ngerman]{babel}
\usepackage{geometry}
\usepackage{fancyhdr}
\usepackage{lastpage}
\usepackage[table]{xcolor}
\usepackage{tcolorbox}
\usepackage{enumitem}
\usepackage{tabularx}
\usepackage{longtable}
\usepackage{array}
\usepackage{booktabs}
\usepackage{hyperref}
\usepackage{ifthen}
\usepackage{amssymb}

% ============================================================================
% KONFIGURATION
% ============================================================================

% --- TIER ---
\newcommand{\plantier}{1}

% --- FACH ---
\newcommand{\fachid}{spanisch}
\newcommand{\fach}{Spanisch}
\newcommand{\fachfarbe}{c41e3a}

% ============================================================================
% VARIABLEN
% ============================================================================

% --- Metadaten ---
\newcommand{\jahrgangsstufe}{10/11/12 (jahrgangsübergreifend)}
\newcommand{\schulart}{Gesamtschule}
\newcommand{\stundenthema}{Cultura y Diversión -- Pufferstunde}
\newcommand{\reihenthema}{Beobachtungsstunde Februar 2026}
\newcommand{\stundennummer}{7}
\newcommand{\reihenstunden}{8}
\newcommand{\unterrichtsdatum}{26.01.2026}
\newcommand{\unterrichtszeit}{[Lt. Stundenplan]}
\newcommand{\raum}{[Fachraum]}

% --- Lehrkraft ---
\newcommand{\lehrkraft}{[Name]}
\newcommand{\ausbildungsstand}{Lehrkraft}

% --- Lerngruppe ---
\newcommand{\lerngruppengroesse}{10}
\newcommand{\maennlich}{1}
\newcommand{\weiblich}{9}
\newcommand{\divers}{0}

% --- Kompetenzen/Niveau ---
\newcommand{\gerniveau}{A2--B1}
\newcommand{\lehrwerk}{A\_tope.com}
\newcommand{\lehrwerkseite}{[entfällt -- kulturelle Stunde]}

% ============================================================================
% LAYOUT
% ============================================================================
\geometry{left=2.5cm, right=2.5cm, top=2.5cm, bottom=2.5cm}
\definecolor{fach}{HTML}{\fachfarbe}
\definecolor{gray}{HTML}{666666}
\definecolor{lightgray}{HTML}{f5f5f5}
\definecolor{successgreen}{HTML}{2a7a4b}
\definecolor{warningorange}{HTML}{b35c00}
\definecolor{gruppeA}{HTML}{2c5aa0}
\definecolor{gruppeB}{HTML}{c41e3a}
\definecolor{puffer}{HTML}{6a0dad}
\definecolor{online}{HTML}{2a7a4b}

\tcbuselibrary{skins,breakable}

% Farbige Boxen
\newtcolorbox{infobox}[1][]{
    colback=lightgray,
    colframe=fach,
    fonttitle=\bfseries,
    title=#1,
    breakable
}

\newtcolorbox{scorebox}{
    colback=white,
    colframe=fach,
    fonttitle=\bfseries,
    title=Didaktik-Score,
    breakable
}

\newtcolorbox{pufferbox}{
    colback=white,
    colframe=puffer,
    fonttitle=\bfseries,
    title={PUFFER -- Entbehrlich für Beobachtungsstunde},
    breakable
}

\newtcolorbox{onlinebox}{
    colback=white,
    colframe=online,
    fonttitle=\bfseries,
    title={ONLINE-VARIANTE -- Bei Lehrerausfall},
    breakable
}

% Kopf-/Fußzeile
\pagestyle{fancy}
\fancyhf{}
\fancyhead[L]{\small\textcolor{gray}{\fach{} -- Jg. \jahrgangsstufe{} -- \stundenthema}}
\fancyhead[R]{\small\textcolor{gray}{Langentwurf Tier 1}}
\fancyfoot[C]{\small\thepage{} / \pageref{LastPage}}
\renewcommand{\headrulewidth}{0.4pt}
\renewcommand{\footrulewidth}{0pt}

% ============================================================================
% DOKUMENT
% ============================================================================
\begin{document}

% ============================================================================
% DECKBLATT
% ============================================================================
\begin{titlepage}
    \centering
    \vspace*{2cm}

    {\large\textcolor{gray}{\schulart}}\\[2cm]

    {\Huge\textcolor{fach}{\textbf{Langentwurf}}}\\[0.5cm]
    {\large Tier 1 -- Vollständig}\\[2cm]

    {\LARGE\textbf{\stundenthema}}\\[0.5cm]
    {\large im Rahmen der Unterrichtsreihe}\\[0.3cm]
    {\Large\textit{\reihenthema}}\\[1cm]

    \textcolor{puffer}{\Large\textbf{--- PUFFER-STUNDE ---}}\\[0.3cm]
    {\normalsize Festigung und Motivation ohne Prüfungsdruck}\\[0.3cm]
    {\small (Entbehrlich für Beobachtungsstunde am 03.02.2026)}\\[1.5cm]

    \begin{tabular}{rl}
        \textbf{Fach:} & \fach{} (\gerniveau) \\[0.3cm]
        \textbf{Klasse:} & \jahrgangsstufe \\[0.3cm]
        \textbf{Lehrwerk:} & \lehrwerk \\[0.3cm]
        \textbf{Datum:} & \underline{\hspace{1cm}\unterrichtsdatum\hspace{1cm}} \\[0.3cm]
        \textbf{Zeit:} & \underline{\hspace{3cm}} (Doppelstunde) \\[0.3cm]
        \textbf{Raum:} & \underline{\hspace{3cm}} \\[0.3cm]
    \end{tabular}

    \vfill

    {\large Lehrkraft: \lehrkraft}\\[0.3cm]
    {\normalsize\textcolor{gray}{\ausbildungsstand}}\\[1cm]

\end{titlepage}

% ============================================================================
% INHALTSVERZEICHNIS
% ============================================================================
\tableofcontents
\newpage

% ============================================================================
% 1. BEDINGUNGSANALYSE
% ============================================================================
\section{Bedingungsanalyse}

\begin{pufferbox}
\textbf{Puffer-Charakter dieser Stunde:}
\begin{itemize}
    \item Diese Stunde ist \textbf{entbehrlich} für die Beobachtungsstunde am 03.02.2026
    \item Der Inhalt baut \textbf{nicht} auf die Beobachtungsstunde auf
    \item Fokus: Festigung, Motivation, kulturelles Lernen, Spaß
    \item Bei Ausfall kann alternativ die Online-Variante eingesetzt werden
\end{itemize}
\end{pufferbox}

\subsection{Lerngruppe}
\begin{tabular}{ll}
    Kursgröße: & \lerngruppengroesse{} SuS (\maennlich{} m / \weiblich{} w / \divers{} d) \\
    Jahrgangsstufe: & Kl. 10, 11, 12 (jahrgangsübergreifend) \\
    Sprachniveau: & \gerniveau{} (GER) -- heterogen \\
    Fremdsprache: & 3. Fremdsprache (Spätbeginner) \\
\end{tabular}

\subsection{Schülerprofile}

\begin{tabular}{|l|c|l|c|p{4.5cm}|}
\hline
\textbf{Name} & \textbf{Kl.} & \textbf{Gruppe} & \textbf{NP} & \textbf{Profil} \\
\hline
\textcolor{gruppeA}{E.H.} & 10 & \textcolor{gruppeA}{A (A2)} & 15 & TOP, Peer-Tutor, Franz. Hintergrund \\
\hline
\textcolor{gruppeA}{L.N.} & 10 & \textcolor{gruppeA}{A (A2)} & 13 & Proaktiv, etwas faul \\
\hline
\textcolor{gruppeA}{C.P.} & 10 & \textcolor{gruppeA}{A (A2)} & 12 & Gute Aussprache, ruhig \\
\hline
\textcolor{gruppeA}{V.S.} & 10 & \textcolor{gruppeA}{A (A2)} & 11 & Konstant, fleißig, ruhig \\
\hline
\textcolor{gruppeB}{L.K.} & 11 & \textcolor{gruppeB}{B (B1)} & 10 & Nachholer, abgelenkt \\
\hline
\textcolor{gruppeB}{H.P.} & 11 & \textcolor{gruppeB}{B (B1)} & 10 & Späteinsteiger, Grundlagenlücken \\
\hline
\textcolor{gruppeB}{A.A.} & 12 & \textcolor{gruppeB}{B (B1)} & 10 & Fleißig, unsicher \\
\hline
\textcolor{gruppeB}{A.S.} & 12 & \textcolor{gruppeB}{B (B1)} & 9 & Einziger Junge, unsicher \\
\hline
\end{tabular}

\vspace{1em}

\textbf{Gruppen-Charakteristik:}
\begin{itemize}
    \item \textcolor{gruppeA}{\textbf{Gruppe A (5 SuS):}} A2-Niveau, Lernjahr 1, starke Aussprache
    \item \textcolor{gruppeB}{\textbf{Gruppe B (5 SuS):}} B1-Niveau, Lernjahr 2--3, mehr Wortschatz aber unsicherer
\end{itemize}

\subsection{Besonderheiten für Pufferstunde}
\begin{itemize}
    \item \textbf{Kein Prüfungsdruck:} Spielerisches Lernen ermöglichen
    \item \textbf{Motivation fördern:} Nach intensiver Arbeit in DS 1--6
    \item \textbf{Gruppenzusammenhalt:} Gemeinsame Aktivitäten für Kohäsion
    \item \textbf{Kulturelles Lernen:} Authentische Einblicke in hispanophone Welt
\end{itemize}

% ============================================================================
% 2. SACHANALYSE
% ============================================================================
\section{Sachanalyse}

\subsection{Kulturelle Inhalte}

\textbf{Musik (Phase 1):}
\begin{itemize}
    \item Spanischsprachige Musik als kultureller Zugang
    \item Empfehlung: \textit{\glqq{}Bailando\grqq{}} (Enrique Iglesias) oder \textit{\glqq{}Vivir mi vida\grqq{}} (Marc Anthony)
    \item Einfache, repetitive Texte mit hohem Wiedererkennungswert
    \item Lückentext als rezeptive Übung (Hörverstehen)
\end{itemize}

\textbf{Geographie (Phase 2):}
\begin{itemize}
    \item 21 spanischsprachige Länder weltweit
    \item Flaggen, Hauptstädte, Kontinente
    \item Kulturelle Vielfalt: Mexiko, Argentinien, Spanien, Kolumbien etc.
    \item Quiz-Format fördert spielerisches Lernen
\end{itemize}

\textbf{Adjektive (Phase 3):}
\begin{itemize}
    \item Wiederholung und Festigung aus DS 6 (Brückenthema)
    \item Bingo als motivierendes Übungsformat
    \item Kombination aus Hören (L ruft) und Erkennen (Adjektiv finden)
\end{itemize}

\textbf{Personenbeschreibung (Phase 4):}
\begin{itemize}
    \item Ratespiel \textit{\glqq{}\textquestiondown{}Quién soy?\grqq{}} (Wer bin ich?)
    \item Bekannte Persönlichkeiten aus spanischsprachigen Ländern
    \item Ja/Nein-Fragen zur Eingrenzung
    \item Anwendung von \textit{ser + Adjektiv} in authentischem Kontext
\end{itemize}

\subsection{Sprachliche Strukturen}

\textbf{Wiederholung (keine Neueinführung):}
\begin{itemize}
    \item Adjektive zur Personenbeschreibung (\textit{simpático, inteligente, famoso...})
    \item \textit{ser + Adjektiv} (Eigenschaften)
    \item Ja/Nein-Fragen: \textit{\textquestiondown{}Es...? \textquestiondown{}Tiene...? \textquestiondown{}Vive en...?}
    \item Länder und Nationalitäten: \textit{español/-a, mexicano/-a, argentino/-a}
\end{itemize}

\subsection{Kulturelle Aspekte}

\begin{tabular}{|l|p{8cm}|}
\hline
\textbf{Thema} & \textbf{Kultureller Hintergrund} \\
\hline
Musik & Reggaetón, Salsa, Pop Latino als globales Phänomen \\
\hline
Geographie & Kolonialgeschichte, Sprachverbreitung, regionale Unterschiede \\
\hline
Persönlichkeiten & Sport (Messi, Nadal), Musik (Shakira), Politik, Kunst \\
\hline
\end{tabular}

% ============================================================================
% 3. DIDAKTISCHE ANALYSE
% ============================================================================
\section{Didaktische Analyse}

\subsection{Stellung in der Sequenz}

Dies ist die \textbf{7. Doppelstunde} der Beobachtungssequenz (8 DS gesamt).

\vspace{0.5em}
\textbf{Sequenzübersicht:}
\begin{center}
\begin{tabular}{|c|l|l|}
\hline
\textbf{DS} & \textbf{Thema} & \textbf{Fokus} \\
\hline
1--5 & Getrennte Progression & A: LJ1 / B: LJ2--3 \\
\hline
6 & Brückenthema: Personas y características & ERSTER OVERLAP \\
\hline
\rowcolor{lightgray} \textbf{7} & \textbf{Cultura y Diversión (PUFFER)} & \textbf{Festigung, Motivation} \\
\hline
8 & Beobachtungsstunde (03.02.2026) & Vollständige Integration \\
\hline
\end{tabular}
\end{center}

\textbf{Funktion dieser Stunde:}
\begin{itemize}
    \item \textbf{Puffer:} Entlastung vor der Beobachtungsstunde
    \item \textbf{Festigung:} Adjektive und Strukturen aus DS 6 festigen
    \item \textbf{Motivation:} Spielerisches Lernen ohne Bewertungsdruck
    \item \textbf{Kulturelles Lernen:} Einblicke in hispanophone Welt
    \item \textbf{Gruppenkohäsion:} Gemeinsame Aktivitäten stärken Zusammenhalt
\end{itemize}

\subsection{Begründung: Festigung ohne Prüfungsdruck}

\textbf{Didaktische Legitimation einer Pufferstunde:}
\begin{enumerate}
    \item \textbf{Konsolidierung:} Nach intensiver Arbeit (DS 1--6) brauchen SuS Zeit zur Festigung
    \item \textbf{Affektive Ziele:} Motivation und Freude am Spanischlernen
    \item \textbf{Kulturelle Kompetenz:} Authentische Einblicke erhöhen intrinsische Motivation
    \item \textbf{Entlastung:} Vor wichtiger Beobachtungsstunde kein zusätzlicher Druck
    \item \textbf{Flexibilität:} Bei Ausfall kann Online-Variante eingesetzt werden
\end{enumerate}

\subsection{Kompetenzziele (reduziert)}

\begin{itemize}
    \item \textbf{Hörverstehen:} Spanisches Lied global verstehen, Lücken ergänzen
    \item \textbf{Landeskunde:} Spanischsprachige Länder kennen und einordnen
    \item \textbf{Sprechen:} Ja/Nein-Fragen stellen und beantworten
    \item \textbf{Wortschatz:} Adjektive zur Personenbeschreibung festigen
    \item \textbf{Interkulturell:} Offenheit für hispanophone Kulturen entwickeln
\end{itemize}

\subsection{Legitimation}

\textbf{Rahmenplan Spanisch MV:}
\begin{itemize}
    \item Interkulturelle Kompetenz: \glqq{}Einblicke in die Lebenswirklichkeit spanischsprachiger Länder\grqq{}
    \item Methodenkompetenz: \glqq{}Spielerische Formen des Lernens nutzen\grqq{}
    \item Affektive Ziele: \glqq{}Interesse an der Zielkultur entwickeln\grqq{}
\end{itemize}

% ============================================================================
% 4. STUNDENZIELE
% ============================================================================
\section{Stundenziele}

\subsection{Grobziel}
Die SuS \textbf{festigen} ihre sprachlichen Kompetenzen und \textbf{erweitern} ihr kulturelles Wissen über die hispanophone Welt durch spielerische und motivierende Aktivitäten.

\subsection{Feinziele (operationalisiert)}

\begin{tabular}{|c|p{7cm}|c|c|}
\hline
\textbf{Nr.} & \textbf{Die SuS können...} & \textbf{AFB} & \textbf{Phase} \\
\hline
FZ1 & ein spanisches Lied global verstehen und Lücken im Text ergänzen & I & 1 (Musik) \\
\hline
FZ2 & mindestens 5 spanischsprachige Länder mit Hauptstädten benennen & I & 2 (Quiz) \\
\hline
FZ3 & Adjektive zur Personenbeschreibung erkennen und zuordnen & I & 3 (Bingo) \\
\hline
FZ4 & gezielte Ja/Nein-Fragen zur Personenidentifikation stellen & II & 4 (Ratespiel) \\
\hline
FZ5 & kulturelle Eigenheiten spanischsprachiger Länder wertschätzend beschreiben & II & alle \\
\hline
\end{tabular}

\vspace{1em}
\textbf{Hinweis:} Die Ziele sind bewusst auf AFB I--II beschränkt (Festigung, keine Progression).

% ============================================================================
% 5. VERLAUFSPLANUNG (PRÄSENZ-VARIANTE)
% ============================================================================
\section{Verlaufsplanung (Präsenz-Variante)}

\textbf{Übersicht Doppelstunde (90 Min.)}

\begin{longtable}{|p{0.8cm}|p{1.5cm}|p{4cm}|p{3cm}|p{1.5cm}|p{2cm}|}
\hline
\textbf{Zeit} & \textbf{Phase} & \textbf{Lehrerhandeln} & \textbf{Schülerhandeln} & \textbf{SF} & \textbf{Medien} \\
\hline
\endhead

\multicolumn{6}{|c|}{\textbf{PHASE 1: MÚSICA EN ESPAÑOL (20 Min.)}} \\
\hline

0--3 & Einstieg & \glqq{}\textexclamdown{}Hola! Hoy tenemos una clase especial: cultura y diversión.\grqq{} Lied ankündigen & Zuhören, Spannung aufbauen & Plenum & -- \\
\hline

3--8 & Hören 1 & Lied abspielen (1x komplett), Aufgabe: \glqq{}Escuchad y disfrutad.\grqq{} & Lied hören, Atmosphäre aufnehmen & Plenum & Audio, Beamer \\
\hline

8--10 & Vorbereitung & Lückentext verteilen, Aufgabe erklären: \glqq{}Completad los huecos.\grqq{} & Lückentext sichten & EA & AB Lückentext \\
\hline

10--15 & Hören 2 & Lied erneut abspielen (ggf. 2x), bei Bedarf pausieren & Lücken ausfüllen & EA & Audio \\
\hline

15--18 & Kontrolle & Lösungen im Plenum besprechen, korrekte Wörter an Tafel & Lösungen vergleichen & Plenum & Tafel \\
\hline

18--20 & Singen & Lied ein letztes Mal abspielen, zum Mitsingen ermutigen & Mitsingen (freiwillig) & Plenum & Audio \\
\hline

\multicolumn{6}{|c|}{\textbf{PHASE 2: PAÍSES HISPANOHABLANTES -- QUIZ-RALLYE (25 Min.)}} \\
\hline

20--23 & Einführung & Weltkarte zeigen: \glqq{}\textquestiondown{}Dónde se habla español?\grqq{} Brainstorming & Länder nennen, L markiert & Plenum & Weltkarte \\
\hline

23--25 & Teams & 2 Teams bilden (gemischt A+B), Quiz-Regeln erklären & Teams finden & Plenum & -- \\
\hline

25--40 & Quiz & Quiz-Rallye durchführen: Flaggen, Hauptstädte, Fakten. Punkte vergeben & Im Team beraten, antworten & Team & Quiz-Karten \\
\hline

40--45 & Auswertung & Punkte zählen, Siegerteam küren, kleine Preise & Jubeln, Reflexion & Plenum & Preise \\
\hline

\multicolumn{6}{|c|}{\textbf{KURZE PAUSE (5 Min.)}} \\
\hline

45--50 & Pause & Pause, Bingo-Karten vorbereiten & Pause & -- & -- \\
\hline

\multicolumn{6}{|c|}{\textbf{PHASE 3: ADJEKTIV-BINGO (25 Min.)}} \\
\hline

50--52 & Vorbereitung & Bingo-Karten verteilen, Regeln erklären & Bingo-Karten sichten & Plenum & Bingo-Karten \\
\hline

52--70 & Bingo & L ruft Adjektive auf Spanisch, SuS markieren. Drei Runden: Reihe, Spalte, Voll & Adjektive erkennen, markieren, \glqq{}\textexclamdown{}Bingo!\grqq{} rufen & EA/Plenum & Bingo-Karten, Stifte \\
\hline

70--75 & Preise & Gewinner erhalten kleine Preise (Bonbons, Sticker) & Preise entgegennehmen & Plenum & Preise \\
\hline

\multicolumn{6}{|c|}{\textbf{PHASE 4: \glqq{}\textquestiondown{}QUIÉN SOY?\grqq{} -- PERSONENRATEN (20 Min.)}} \\
\hline

75--78 & Einführung & Spiel erklären: Berühmte Person auf Stirn, nur Ja/Nein-Fragen & Regeln verstehen & Plenum & -- \\
\hline

78--80 & Vorbereitung & Kärtchen mit Persönlichkeiten verteilen (verdeckt), Paare/Gruppen bilden & Kärtchen nehmen & Plenum & Personenkärtchen \\
\hline

80--88 & Spielen & SuS spielen in Kleingruppen, L beobachtet und hilft bei Fragen & Fragen stellen: \glqq{}\textquestiondown{}Es hombre?\grqq{} \glqq{}\textquestiondown{}Es famoso?\grqq{} etc. & GA & Kärtchen \\
\hline

88--90 & Abschluss & Zusammenfassung: \glqq{}\textquestiondown{}Qué habéis aprendido hoy?\grqq{} Ausblick DS 8 & Reflektieren & Plenum & -- \\
\hline

\end{longtable}

\textbf{Legende:} EA = Einzelarbeit, GA = Gruppenarbeit, SF = Sozialform

% ============================================================================
% 6. ONLINE-VARIANTE (BEI LEHRERAUSFALL)
% ============================================================================
\section{Online-Variante (bei Lehrerausfall)}

\begin{onlinebox}
\textbf{Einsatzszenario:} Bei Lehrerausfall können die SuS diese Stunde selbstständig online bearbeiten. Die Aufgaben werden über itslearning bereitgestellt.

\vspace{1em}
\textbf{Gesamtdauer:} ca. 65--70 Minuten (ohne Abgabe)
\end{onlinebox}

\subsection{Aufgabenübersicht}

\begin{tabular}{|c|l|c|p{6cm}|}
\hline
\textbf{Nr.} & \textbf{Aufgabe} & \textbf{Zeit} & \textbf{Beschreibung} \\
\hline
1 & LearningApps & 20' & 3 Zuordnungsübungen (Länder, Flaggen, Adjektive) \\
\hline
2 & Quizlet & 20' & Vokabel-Review: Adjektive + Länder \\
\hline
3 & Video + Fragen & 25' & Kurzfilm + 5 Verständnisfragen \\
\hline
\multicolumn{2}{|r|}{\textbf{Gesamt:}} & \textbf{65'} & \\
\hline
\end{tabular}

\subsection{Detailbeschreibung}

\textbf{Aufgabe 1: LearningApps (20 Min.)}
\begin{enumerate}
    \item \textbf{Zuordnung Flaggen:} 10 Flaggen den Ländern zuordnen
    \item \textbf{Zuordnung Hauptstädte:} 10 Länder-Hauptstadt-Paare
    \item \textbf{Adjektiv-Memory:} Spanisch-Deutsch Paare finden
\end{enumerate}
\textit{Links werden in itslearning bereitgestellt.}

\vspace{0.5em}
\textbf{Aufgabe 2: Quizlet (20 Min.)}
\begin{itemize}
    \item Lernmodus: \glqq{}Lernen\grqq{} (adaptive Karteikarten)
    \item Set 1: Adjektive zur Personenbeschreibung (20 Wörter)
    \item Set 2: Países hispanohablantes (15 Länder + Hauptstädte)
    \item Abschluss: \glqq{}Test\grqq{}-Modus (Screenshot für Abgabe)
\end{itemize}

\vspace{0.5em}
\textbf{Aufgabe 3: Video + Verständnisfragen (25 Min.)}
\begin{itemize}
    \item Video: \textit{\glqq{}Un día en la vida de...\grqq{}} (ca. 8 Min., Spanisch mit Untertiteln)
    \item Alternative: \textit{\glqq{}Destino: [Land]\grqq{}} (Reisedoku, 10 Min.)
\end{itemize}

\textbf{Verständnisfragen (auf Deutsch zu beantworten):}
\begin{enumerate}
    \item In welchem Land spielt das Video?
    \item Welche Sehenswürdigkeiten werden gezeigt?
    \item Welche spanischen Wörter hast du verstanden? (Mind. 5)
    \item Was hat dich überrascht?
    \item Würdest du dieses Land besuchen wollen? Warum (nicht)?
\end{enumerate}

\subsection{Abgabe via itslearning}

\begin{itemize}
    \item \textbf{Frist:} Bis 18:00 Uhr am selben Tag
    \item \textbf{Format:} Kurze Nachricht mit Screenshots (LearningApps, Quizlet) + Antworten auf Videofragen
    \item \textbf{Bewertung:} Keine Note, aber Rückmeldung bei Nicht-Abgabe
\end{itemize}

% ============================================================================
% 7. ERWARTETE SCHWIERIGKEITEN
% ============================================================================
\section{Erwartete Schwierigkeiten}

\begin{tabular}{|p{4.5cm}|p{8.5cm}|}
\hline
\textbf{Schwierigkeit} & \textbf{Reaktion} \\
\hline
\textbf{SuS nehmen Pufferstunde nicht ernst} & Klare Struktur beibehalten, Preise als Anreiz, Beteiligung loben \\
\hline
\textbf{Lied zu schnell/schwer} & Langsames Lied wählen, mehrfach abspielen, Pause-Taste nutzen \\
\hline
\textbf{Quiz-Fragen zu schwer} & Verschiedene Schwierigkeitsgrade vorbereiten (leicht/mittel/schwer) \\
\hline
\textbf{Bingo geht zu schnell} & Adjektive langsam und deutlich rufen, ggf. an Tafel schreiben \\
\hline
\textbf{Ratespiel stockt} & Hilfestellung: \glqq{}Pregunta: \textquestiondown{}Es de España?\grqq{} Beispielfragen an Tafel \\
\hline
\textbf{Online-Variante: SuS arbeiten nicht} & Abgabepflicht mit Screenshot-Nachweis, Elterninformation bei Nicht-Abgabe \\
\hline
\textbf{Technische Probleme (Online)} & Alternative: AB per E-Mail/PDF mit analogen Aufgaben \\
\hline
\end{tabular}

% ============================================================================
% 8. HAUSAUFGABE
% ============================================================================
\section{Hausaufgabe}

\textbf{Keine formale Hausaufgabe} (Pufferstunde).

\vspace{0.5em}
\textbf{Optional/freiwillig:}
\begin{itemize}
    \item Eigenes Lieblingslied auf Spanisch suchen und mitbringen
    \item 3 interessante Fakten über ein spanischsprachiges Land recherchieren
\end{itemize}

% ============================================================================
% 9. ANLAGEN
% ============================================================================
\section{Anlagen}

\begin{enumerate}
    \item Lückentext Lied (mit Lösungen)
    \item Quiz-Karten: Países hispanohablantes
    \item Bingo-Karten (6 verschiedene, je 2x drucken)
    \item Personenkärtchen für \glqq{}\textquestiondown{}Quién soy?\grqq{}
    \item Online-Aufgaben (Links für itslearning)
\end{enumerate}

\subsection{Lückentext Lied (Beispiel: \glqq{}Vivir mi vida\grqq{})}

\begin{verbatim}
VIVIR MI VIDA (Marc Anthony)
=============================

Voy a _________ (1), voy a bailar
Vivir mi _________ (2), la la la la
Voy a _________ (3), voy a gozar
Vivir mi _________ (4), la la la la

A veces llega la _________ (5)
A veces me _________ (6) yo también
Mas yo sé bien
Que hay mañana

Para qué _________ (7), para qué sufrir
Empieza a _________ (8), empieza a _________ (9)
Vive, _________ (10), ríe, ama...

LÖSUNGEN:
(1) reír, (2) vida, (3) bailar, (4) vida, (5) lluvia,
(6) siento triste, (7) llorar, (8) soñar, (9) sentir, (10) sueña
\end{verbatim}

\subsection{Quiz-Karten: Países hispanohablantes}

\textbf{Beispielfragen (unterschiedliche Schwierigkeitsgrade):}

\begin{verbatim}
LEICHT (1 Punkt):
- ¿Cuál es la capital de España? [Madrid]
- ¿En qué continente está Argentina? [Sudamérica]
- ¿De qué país es esta bandera? [Zeige Mexiko-Flagge]

MITTEL (2 Punkte):
- ¿Cuál es la capital de Perú? [Lima]
- ¿Qué países tienen frontera con España? [Portugal, Francia, Andorra]
- ¿En qué país está Machu Picchu? [Perú]

SCHWER (3 Punkte):
- ¿Cuál es la capital de Honduras? [Tegucigalpa]
- ¿Cuántos países hispanohablantes hay en total? [21]
- ¿Qué país hispanohablante tiene el español como lengua oficial
  pero no es su lengua mayoritaria? [Guinea Ecuatorial]
\end{verbatim}

\subsection{Bingo-Karten Adjektive}

\textbf{Wortpool (24 Adjektive):}
\begin{verbatim}
simpático/-a    inteligente     trabajador/-a    paciente
amable          divertido/-a    creativo/-a      responsable
organizado/-a   serio/-a        tranquilo/-a     reservado/-a
perezoso/-a     impaciente      egoísta          generoso/-a
famoso/-a       rico/-a         guapo/-a         feo/-a
alto/-a         bajo/-a         joven            viejo/-a
\end{verbatim}

\textbf{Bingo-Karte 1 (5x5):}
\begin{verbatim}
+-------------+-------------+-------------+-------------+-------------+
| simpático   | inteligente | trabajador  | GRATIS      | amable      |
+-------------+-------------+-------------+-------------+-------------+
| divertido   | creativo    | responsable | organizado  | serio       |
+-------------+-------------+-------------+-------------+-------------+
| tranquilo   | reservado   | GRATIS      | perezoso    | impaciente  |
+-------------+-------------+-------------+-------------+-------------+
| egoísta     | generoso    | famoso      | rico        | guapo       |
+-------------+-------------+-------------+-------------+-------------+
| alto        | bajo        | joven       | viejo       | paciente    |
+-------------+-------------+-------------+-------------+-------------+
\end{verbatim}

\textit{(5 weitere Karten mit unterschiedlicher Anordnung erstellen)}

\subsection{Personenkärtchen \glqq{}\textquestiondown{}Quién soy?\grqq{}}

\begin{verbatim}
+------------------------+    +------------------------+
| LIONEL MESSI           |    | SHAKIRA                |
| Fútbol, Argentina      |    | Música, Colombia       |
+------------------------+    +------------------------+

+------------------------+    +------------------------+
| RAFAEL NADAL           |    | PABLO PICASSO          |
| Tenis, España          |    | Arte, España           |
+------------------------+    +------------------------+

+------------------------+    +------------------------+
| FRIDA KAHLO            |    | JULIO IGLESIAS         |
| Arte, México           |    | Música, España         |
+------------------------+    +------------------------+

+------------------------+    +------------------------+
| PENÉLOPE CRUZ          |    | GABRIEL GARCÍA MÁRQUEZ |
| Cine, España           |    | Literatura, Colombia   |
+------------------------+    +------------------------+
\end{verbatim}

\textbf{Beispielfragen für das Spiel:}
\begin{itemize}
    \item \textit{\textquestiondown{}Soy hombre?} (Bin ich ein Mann?)
    \item \textit{\textquestiondown{}Soy de España?} (Bin ich aus Spanien?)
    \item \textit{\textquestiondown{}Soy deportista?} (Bin ich Sportler?)
    \item \textit{\textquestiondown{}Estoy vivo/-a?} (Lebe ich noch?)
    \item \textit{\textquestiondown{}Soy cantante?} (Bin ich Sänger?)
\end{itemize}

% ============================================================================
% 10. DIDAKTIK-SCORE
% ============================================================================
\newpage
\section{Didaktik-Score}

\begin{scorebox}
\textbf{Bewertung der didaktischen Qualität nach 10 Dimensionen}\\
\textbf{Ziel-Score: $\geq$ 9.0 (Premium-Qualität)}

\vspace{1em}
\begin{tabular}{|c|l|c|c|p{4.5cm}|}
\hline
\textbf{\#} & \textbf{Dimension} & \textbf{Gew.} & \textbf{Score} & \textbf{Notiz} \\
\hline
1 & Differenzierung & 15\% & 8/10 & Gemeinsame Aktivitäten, aber Quiz mit Schwierigkeitsgraden \\
\hline
2 & Taxonomie (AFB) & 10\% & 8/10 & Bewusst I--II (Festigung), kein III \\
\hline
3 & Lernziel-Alignment & 10\% & 10/10 & Passt zu RP: Kulturelle Kompetenz \\
\hline
4 & Aktivierung & 10\% & 10/10 & Alle Phasen aktiv, spielerisch \\
\hline
5 & Scaffolding & 10\% & 9/10 & Hilfestellungen bei allen Spielen \\
\hline
6 & Feedback-Qualität & 10\% & 8/10 & Sofortiges Feedback durch Spielformat \\
\hline
7 & Sprachliche Klarheit & 10\% & 10/10 & Einfache Anweisungen, spielerisch \\
\hline
8 & Motivation \& Relevanz & 10\% & 10/10 & Hohe Motivation durch Spiele, Preise, Musik \\
\hline
9 & Inklusion (UDL) & 10\% & 9/10 & Auditiv + visuell + kinästhetisch \\
\hline
10 & Praktikabilität & 5\% & 10/10 & Zeitrahmen großzügig, Material einfach \\
\hline
\multicolumn{3}{|r|}{\textbf{GESAMT:}} & \textbf{9.1/10} & \\
\hline
\end{tabular}

\vspace{1em}
\textbf{Bewertungsschwellen:}
\begin{itemize}
    \item[\checkmark] \textcolor{successgreen}{$\geq$ 9.0: \textbf{FREIGABE} (Premium-Qualität)}
    \item[$\Box$] \textcolor{warningorange}{8.0--8.9: Iteration empfohlen}
    \item[$\Box$] 6.0--7.9: Überarbeitung erforderlich
    \item[$\Box$] $<$ 6.0: Grundlegende Neukonzeption
\end{itemize}

\vspace{1em}
\textcolor{successgreen}{\textbf{STATUS: FREIGABE}}

\vspace{1em}
\textbf{Stärken:}
\begin{itemize}
    \item Hohe Motivation durch spielerische Formate
    \item Kulturelles Lernen authentisch eingebunden
    \item Online-Variante für Lehrerausfall vorbereitet
    \item Klare Puffer-Funktion ohne Prüfungsdruck
    \item Festigung der DS-6-Inhalte (Adjektive) integriert
\end{itemize}

\textbf{Hinweis zum Score:}
\begin{itemize}
    \item Differenzierung bewusst reduziert (gemeinsame Aktivitäten)
    \item AFB III entfällt planmäßig (Puffer = Festigung)
    \item Für eine Pufferstunde ist dieser Score angemessen
\end{itemize}

\end{scorebox}

\end{document}
